\section{根号を含む積分}\label{節:平方根と部分積分}

$\sqrt{1+x^2}$ も,そのままでは基本公式の形ではない.前章の三角置換から $x=\tan\theta$ を考えることはできるが,ここでは部分積分を使う別の道を調べよう.一般的な形は,\,$A>0$ として\footnote{$A<0$ のときは $x+\sqrt{x^2+A}$ が負になる範囲があり,対数の中身に絶対値が必要になる.本節では $A>0$ に限る.}
\[
 \int\sqrt{x^2+A}\dd x
\]
で与えられる.循環する積分と同じく,部分積分でもとの積分を含む等式を作り,それを解くことを目標とする.

手がかりは,根号を微分すると $\frac{x}{\sqrt{x^2+A}}$ になり,さらに $x$ を掛ければ分子に $x^2$ ができることである.この $x^2$ を $(x^2+A)-A$ と分けると,もとの根号が再び現れる.

$u=\sqrt{x^2+A}$,\,$v'=1$ と選ぶ.\,$u'=\frac{x}{\sqrt{x^2+A}}$,\,$v=x$ だから,部分積分の公式より
\[
 \int\sqrt{x^2+A}\dd x
 =x\sqrt{x^2+A}-\int\frac{x^2}{\sqrt{x^2+A}}\dd x
\]
を得る.右辺の積分を見ただけでは簡単になったか分かりにくいが,分子を分けると
\[
\begin{aligned}
 \frac{x^2}{\sqrt{x^2+A}}
 &=\frac{(x^2+A)-A}{\sqrt{x^2+A}}\\
 &=\sqrt{x^2+A}-\frac{A}{\sqrt{x^2+A}}
\end{aligned}
\]
となる.したがって,もとの積分を左辺へ移項すれば,任意定数を除いて
\[
 2\int\sqrt{x^2+A}\dd x
 =x\sqrt{x^2+A}+A\int\frac1{\sqrt{x^2+A}}\dd x
\]
という関係が得られる.未知の積分がもう一度現れたが,その係数をまとめることができた.残る課題は,根号の逆数の積分である.

この原始関数は少々発想しにくいので,結果からお見せしよう.
\[
 \int\frac1{\sqrt{x^2+A}}\dd x
 =\log\left(x+\sqrt{x^2+A}\right)+C
\]
である.実際,\,$r(x)=\sqrt{x^2+A}$ と略記すると $r'(x)=\frac{x}{r(x)}$ だから,
\[
\begin{aligned}
 \frac{\dd}{\dd x}\log\{x+r(x)\}
 &=\cfrac{1+\cfrac{x}{r(x)}}{x+r(x)}\\
 &=\frac{1}{x+r(x)}\cdot\frac{r(x)+x}{r(x)}\\
 &=\frac1{r(x)}
\end{aligned}
\]
となる.また $r(x)=\sqrt{x^2+A}>|x|$ なので,\,$x+r(x)>0$ はすべての実数 $x$ で成り立ち,対数の中身は正である.これを先ほどの関係へ戻すと,
\[
 \int\sqrt{x^2+A}\dd x
 =\frac12\left\{x\sqrt{x^2+A}
       +A\log\left(x+\sqrt{x^2+A}\right)\right\}+C
\]
を得る.

\begin{bookpoint}
\textbf{POINT}\par
根号を $u$,\,$1$ を $v'$ として部分積分し,分子を $x^2=(x^2+A)-A$ と分けてもとの積分を移項する.残る根号の逆数は $\log\left(x+\sqrt{x^2+A}\right)$ で積分する.
\end{bookpoint}

\begin{example}{根号の部分積分}
$\int\sqrt{1+x^2}\dd x$ を求めよ.
\end{example}

\paragraphbf{方針}
$\sqrt{1+x^2}$ を微分し,\,$1$ を積分する.分子を $x^2=(1+x^2)-1$ と分け,もとの積分を移項する.

\paragraphbf{解答}
$\frac12\left\{x\sqrt{1+x^2}+\log\left(x+\sqrt{1+x^2}\right)\right\}+C$

\paragraphbf{解法}
$u=\sqrt{1+x^2}$,\,$v'=1$ と選ぶ.\,$u'=\frac{x}{\sqrt{1+x^2}}$,\,$v=x$ だから,
\[
 \int\sqrt{1+x^2}\dd x
 =x\sqrt{1+x^2}-\int\frac{x^2}{\sqrt{1+x^2}}\dd x
\]
である.分子を分けると
\[
 \frac{x^2}{\sqrt{1+x^2}}
 =\sqrt{1+x^2}-\frac1{\sqrt{1+x^2}}
\]
となるので,もとの積分を移項して
\[
 2\int\sqrt{1+x^2}\dd x
 =x\sqrt{1+x^2}+\int\frac1{\sqrt{1+x^2}}\dd x
\]
を得る.右辺の積分は $\log\left(x+\sqrt{1+x^2}\right)+C$ であり,\,$x+\sqrt{1+x^2}>0$ がすべての実数 $x$ で成り立つ.よって
\[
 \int\sqrt{1+x^2}\dd x
 =\frac12\left\{x\sqrt{1+x^2}
       +\log\left(x+\sqrt{1+x^2}\right)\right\}+C
\]
である.最後に答え全体を微分すると,
\[
\begin{aligned}
 &\frac12\left\{\sqrt{1+x^2}
       +\frac{x^2}{\sqrt{1+x^2}}
       +\frac1{\sqrt{1+x^2}}\right\}\\
 &\qquad=\frac12\{\sqrt{1+x^2}+\sqrt{1+x^2}\}\\
 &\qquad=\sqrt{1+x^2}
\end{aligned}
\]
となる.この確認により,途中で不定積分を移項した計算も,最終的な原始関数として正しいことが分かる.

\paragraphbf{対数の式の発想}
$\log\left(x+\sqrt{x^2+A}\right)$ という式は,どこから出てきたのだろうか.\,$\frac1{r(x)}$ を積分したいので,微分すると分母に $r(x)$ が現れる式を探す.\,$r'(x)=\frac{x}{r(x)}$ はその候補だが,分子に $x$ が残る.そこで $x'=1$ も,分母を $r(x)$ にして $\frac{r(x)}{r(x)}$ と並べてみよう.この2つを足せば,
\[
 \{x+r(x)\}'=\frac{x+r(x)}{r(x)}
\]
となり,微分する前の $x+r(x)$ が分子に現れる.ここで対数の微分 $\{\log g(x)\}'=\frac{g'(x)}{g(x)}$ を思い出せば,\,$g(x)=x+r(x)$ として分子と分母の $x+r(x)$ が約分され,\,$\frac1{r(x)}$ だけが残るのである.

\paragraphbf{他の置換との関係}
同じ積分は置換でも求められる.\,$x=\tan\theta$ と置くと $\int\sqrt{1+x^2}\dd x$ は $\int\frac{\dd\theta}{\cos^3\theta}$ となり,ここでも部分積分が必要になる.また,本節で現れた $x+\sqrt{1+x^2}$ そのものを $t$ と置く置換もある.これらは\chapterref{章:置換の背景}で双曲線関数とともに扱い,本節の答えと比較する.
